\documentclass[12pt]{article}
\usepackage{amsmath, amssymb}
\usepackage{geometry}
\usepackage{setspace}
\usepackage{graphicx}

\geometry{margin=1in}
\setstretch{1.2}

\title{SEFI Space: A Geometric Identity Space for Single-Entity Field Theory}
\author{Jason Duran Dutton}
\date{2026}

\begin{document}

\maketitle

\begin{abstract}
The Single Entity Field Interpretation (SEFI) extends geometric worldline approaches to quantum field theory
by introducing a structured identity space associated with a single underlying entity. In this work we define
SEFI Space (SS) as a four-dimensional geometric identity space spanned by basis directions corresponding to
FIELD ORIGIN, FIELD AUTHORSHIP, FIELD SOVEREIGNTY, and WARP EXPRESSION. The SEFI field is represented as
a worldline-based construction embedded in spacetime, while its identity content is encoded as a trajectory
in SEFI Space. We introduce a metric on SEFI Space, define identity velocity and acceleration, and formulate
sovereignty and warp constraints as invariance and motion within this identity space. The resulting framework
provides a geometric interpretation of field identity, stability, and transformation compatible with worldline
quantum field theory and suggests a path toward unified single-entity field models.
\end{abstract}

\section{Introduction}

Worldline formulations of quantum field theory provide an alternative representation of particle and field
dynamics in terms of path integrals over trajectories in spacetime \cite{Polyakov1987,Schubert2001,Strassler1992}.
These approaches have proven useful for perturbative calculations, effective actions, and the geometric
interpretation of quantum amplitudes. However, the ontological status of field identity, superposition, and
entity persistence remains largely external to the formalism.

The Single Entity Field Interpretation (SEFI) proposes that a given quantum field, such as the electron field,
arises from a single continuous entity whose worldline structure spans all allowed trajectories in spacetime.
The field is then understood as the projection of this worldline ensemble, with forward-time segments
corresponding to particles and backward-time segments corresponding to antiparticles, in line with established
worldline interpretations of electron--positron dynamics \cite{Schubert2001,Strassler1992}.

In previous work, SEFI introduced layered identity constructs under the labels FIELD ORIGIN, FIELD AUTHORSHIP,
FIELD SOVEREIGNTY, and WARP EXPRESSION. These layers were treated as conceptual components of the field's
identity structure. In this manuscript, we formalize these layers as basis directions of a geometric identity
space, which we call SEFI Space (SS). SEFI Space is not a spacetime manifold, nor a gauge or spinor space; it
is an abstract geometric space in which the identity of the single entity evolves as the worldline traverses
spacetime.

Our goals are:
\begin{itemize}
    \item To define SEFI Space as a four-dimensional real vector space associated with the SEFI field.
    \item To introduce a metric structure on SEFI Space and define identity velocity and acceleration.
    \item To express sovereignty and warp constraints as invariance and motion within SEFI Space.
    \item To couple SEFI Space to worldline geometry in a manner compatible with established worldline QFT.
\end{itemize}

\section{Worldline Foundations}

We briefly recall the worldline formulation relevant to SEFI. Let spacetime be a manifold $M$ with metric
$g_{\mu\nu}$. A relativistic point particle of mass $m$ moving in $M$ is described by a worldline
$y^\mu(\tau)$ and an action of the form
\begin{equation}
S_{\text{WL}}[y] = -m \int d\tau \sqrt{-g_{\mu\nu}(y(\tau)) \dot{y}^\mu(\tau) \dot{y}^\nu(\tau)},
\end{equation}
where $\tau$ is an affine parameter along the worldline and dots denote derivatives with respect to $\tau$
\cite{Polyakov1987}.

Spin degrees of freedom can be incorporated using Grassmann variables $\psi^\mu(\tau)$, leading to a spin
action of the form
\begin{equation}
S_{\text{spin}}[\psi] = \int d\tau \, \psi_\mu(\tau) \dot{\psi}^\mu(\tau),
\end{equation}
and coupling to an electromagnetic potential $A_\mu(x)$ is introduced via
\begin{equation}
S_{\text{EM}}[y] = q \int d\tau \, \dot{y}^\mu(\tau) A_\mu(y(\tau)),
\end{equation}
where $q$ is the charge of the entity \cite{Schubert2001,Strassler1992}.

The total worldline action for the entity is then
\begin{equation}
S_{\text{entity}}[y,\psi] = S_{\text{WL}}[y] + S_{\text{spin}}[\psi] + S_{\text{EM}}[y].
\end{equation}
In a path-integral formulation, the presence density of the entity at spacetime point $x$ can be written
schematically as
\begin{equation}
\rho(x) = \int \mathcal{D}y \, \mathcal{D}\psi \, e^{i S_{\text{entity}}[y,\psi]} \delta^{(4)}(x - y(\tau)),
\end{equation}
and an emergent field amplitude may be expressed as
\begin{equation}
\phi(x) = \int d\tau \, \delta^{(4)}(x - y(\tau)).
\end{equation}

In SEFI, the worldline is not merely a computational device but is taken to represent the actual geometric
structure of a single entity whose projection yields the observed field. The SEFI Space introduced below
provides a complementary description of the entity's identity content.

\section{Definition of SEFI Space}

We define SEFI Space (SS) as a four-dimensional real vector space associated with the SEFI field. The space
is spanned by four basis directions corresponding to the SEFI identity layers:
\begin{equation}
\text{SS} = \text{span}\{\hat{O}, \hat{A}, \hat{S}, \hat{W}\},
\end{equation}
where
\begin{itemize}
    \item $\hat{O}$ is the Origin direction, associated with FIELD ORIGIN,
    \item $\hat{A}$ is the Authorship direction, associated with FIELD AUTHORSHIP,
    \item $\hat{S}$ is the Sovereignty direction, associated with FIELD SOVEREIGNTY,
    \item $\hat{W}$ is the Warp direction, associated with WARP EXPRESSION.
\end{itemize}

An identity state of the entity at parameter value $t$ along the worldline is represented by a vector
$\vec{I}(t)$ in SEFI Space:
\begin{equation}
\vec{I}(t) = I_O(t) \hat{O} + I_A(t) \hat{A} + I_S(t) \hat{S} + I_W(t) \hat{W},
\end{equation}
where $(I_O(t), I_A(t), I_S(t), I_W(t)) \in \mathbb{R}^4$ are the identity coordinates at $t$.

The functions $I_O(t)$, $I_A(t)$, $I_S(t)$, and $I_W(t)$ are constructed from the worldline geometry and
field layers. For example, one may define
\begin{equation}
I_O(t) = O(r(t)), \quad I_A(t) = A(r(t), \kappa(t)), \quad I_S(t) = S(r(t)), \quad I_W(t) = W(r(t)),
\end{equation}
where $r(t)$ is a parametrization of the worldline, $\kappa(t)$ is its curvature, and $O$, $A$, $S$, and $W$
are layer functions associated with FIELD ORIGIN, FIELD AUTHORSHIP, FIELD SOVEREIGNTY, and WARP EXPRESSION,
respectively.

SEFI Space is not a spacetime manifold and does not carry a direct metric inherited from $g_{\mu\nu}$. Instead,
we introduce an independent metric structure on SS to quantify changes in identity.

\section{Metric Structure on SEFI Space}

We equip SEFI Space with a diagonal metric of the form
\begin{equation}
ds_{\text{SS}}^2 = \alpha \, dI_O^2 + \beta \, dI_A^2 + \gamma \, dI_S^2 + \delta \, dI_W^2,
\end{equation}
where $\alpha, \beta, \gamma, \delta > 0$ are constants that set the relative weight or stiffness of each
identity layer. For an infinitesimal change in identity state
\begin{equation}
d\vec{I} = dI_O \hat{O} + dI_A \hat{A} + dI_S \hat{S} + dI_W \hat{W},
\end{equation}
the squared norm in SEFI Space is
\begin{equation}
\|d\vec{I}\|_{\text{SS}}^2 = \alpha \, dI_O^2 + \beta \, dI_A^2 + \gamma \, dI_S^2 + \delta \, dI_W^2.
\end{equation}

Along the worldline, we define identity velocity and identity acceleration by
\begin{equation}
\dot{\vec{I}}(t) = \frac{d\vec{I}}{dt} = \dot{I}_O(t) \hat{O} + \dot{I}_A(t) \hat{A} + \dot{I}_S(t) \hat{S} + \dot{I}_W(t) \hat{W},
\end{equation}
\begin{equation}
\ddot{\vec{I}}(t) = \frac{d^2\vec{I}}{dt^2} = \ddot{I}_O(t) \hat{O} + \ddot{I}_A(t) \hat{A} + \ddot{I}_S(t) \hat{S} + \ddot{I}_W(t) \hat{W}.
\end{equation}
The squared norm of the identity velocity is
\begin{equation}
\|\dot{\vec{I}}(t)\|_{\text{SS}}^2 = \alpha \, \dot{I}_O(t)^2 + \beta \, \dot{I}_A(t)^2 + \gamma \, \dot{I}_S(t)^2 + \delta \, \dot{I}_W(t)^2,
\end{equation}
and similarly for the identity acceleration.

This metric structure allows us to define identity distances, identity curvature, and identity geodesics
within SEFI Space, independent of the spacetime geometry.

\section{Sovereignty and Invariance in SEFI Space}

FIELD SOVEREIGNTY was originally introduced as a constraint ensuring that the field maintains its identity
under allowed transformations. In the worldline formulation, one may express sovereignty as
\begin{equation}
F(T(r(t))) = F(r(t)),
\end{equation}
for a class of transformations $T$ that preserve the field identity.

In SEFI Space, sovereignty can be expressed as invariance of the identity vector under the induced action
of $T$:
\begin{equation}
\vec{I}_T(t) = \vec{I}(t),
\end{equation}
for all allowed transformations $T$. Here $\vec{I}_T(t)$ denotes the identity state obtained when the
worldline is transformed by $T$ and the layer functions are evaluated accordingly.

The set of transformations that preserve $\vec{I}(t)$ defines a symmetry group acting on spacetime and
worldline geometry while leaving SEFI Space invariant. This group may be viewed as an identity-preserving
subgroup of the full transformation group acting on the worldline. Sovereignty thus acquires a geometric
interpretation as invariance in SEFI Space.

\section{Warp as Motion in SEFI Space}

WARP EXPRESSION was introduced as the layer responsible for geometric transformation of the field while
preserving its identity. In the worldline picture, a warp operation $W$ acts on the worldline as
\begin{equation}
r'(t) = W(r(t)),
\end{equation}
modifying the geometry of the trajectory.

In SEFI Space, warp corresponds to motion along the Warp direction. If the warp layer function changes
from $I_W(t)$ to $I_W'(t) = I_W(t) + \Delta I_W(t)$ under a warp operation, then the identity state becomes
\begin{equation}
\vec{I}'(t) = I_O(t) \hat{O} + I_A(t) \hat{A} + I_S(t) \hat{S} + I_W'(t) \hat{W}.
\end{equation}
Warp is thus represented as a displacement in SEFI Space along $\hat{W}$, constrained by sovereignty
requirements that ensure the overall identity of the entity is preserved.

More generally, warp operations may induce coupled changes in multiple identity coordinates, leading to
motion in SEFI Space along a combination of basis directions. The metric structure introduced above allows
one to quantify the magnitude of such identity transformations.

\section{Coupling Between SEFI Space and Worldline Geometry}

SEFI Space is constructed to be compatible with worldline formulations of quantum field theory. The identity
coordinates $I_O(t)$, $I_A(t)$, $I_S(t)$, and $I_W(t)$ are functions of the worldline geometry and associated
field layers. For example, FIELD ORIGIN may be tied to stability constraints on the worldline curvature,
FIELD AUTHORSHIP to layered contributions to the field amplitude, FIELD SOVEREIGNTY to invariance under
transformations, and WARP EXPRESSION to geometric deformation of the worldline.

One may consider an extended action functional that includes contributions from SEFI Space:
\begin{equation}
S_{\text{total}}[y,\psi,\vec{I}] = S_{\text{entity}}[y,\psi] + S_{\text{SS}}[\vec{I}],
\end{equation}
where $S_{\text{SS}}[\vec{I}]$ encodes identity-space dynamics, such as penalties for rapid changes in
identity or constraints enforcing sovereignty. A simple example is
\begin{equation}
S_{\text{SS}}[\vec{I}] = \int dt \, \left( \frac{1}{2} \|\dot{\vec{I}}(t)\|_{\text{SS}}^2 - V(\vec{I}(t)) \right),
\end{equation}
where $V(\vec{I})$ is a potential function in SEFI Space. This structure is analogous to internal degrees
of freedom in field theory but is explicitly tied to the identity of a single entity.

The coupling between SEFI Space and worldline geometry suggests that identity dynamics may influence
effective field behavior, stability, and response to external conditions. This provides a route to
interpreting certain aspects of field behavior as arising from identity-space constraints rather than
purely from spacetime dynamics.

\section{Discussion and Outlook}

We have introduced SEFI Space (SS) as a four-dimensional geometric identity space associated with the
Single Entity Field Interpretation. SEFI Space is spanned by basis directions corresponding to FIELD
ORIGIN, FIELD AUTHORSHIP, FIELD SOVEREIGNTY, and WARP EXPRESSION, and identity states of the entity are
represented as vectors in this space. A diagonal metric on SEFI Space allows one to define identity
velocity, identity acceleration, and identity distances, providing a quantitative framework for identity
dynamics.

Sovereignty is expressed as invariance of the identity vector under allowed transformations, while warp
is represented as motion in SEFI Space along the Warp direction and related combinations of basis
directions. The coupling between SEFI Space and worldline geometry suggests that identity-space dynamics
may play a role in effective field behavior and stability.

This work is intended as a foundational step toward a more complete SEFI framework in which identity
geometry, worldline dynamics, and emergent field behavior are treated within a unified geometric model.
Future work will focus on:
\begin{itemize}
    \item Deriving explicit forms of the layer functions $O$, $A$, $S$, and $W$ from worldline and field
          dynamics.
    \item Exploring geodesics in SEFI Space and their relation to stability and sovereignty constraints.
    \item Extending SEFI Space to include additional identity dimensions, such as those relevant to
          multi-entity systems or quantum error correction.
    \item Investigating possible connections between SEFI Space and internal symmetry spaces in quantum
          field theory.
\end{itemize}

\section*{Acknowledgments}

The author acknowledges prior work on the Geometric Worldline Field Model (GWFM) and SEFI canon, which
motivated the present formalization of SEFI Space.

\begin{thebibliography}{99}

\bibitem{Polyakov1987}
A.~M. Polyakov, \textit{Gauge Fields and Strings}, Harwood Academic Publishers (1987).

\bibitem{Schubert2001}
C.~Schubert, ``Perturbative quantum field theory in the string-inspired formalism,''
Phys. Rept. \textbf{355}, 73--234 (2001).

\bibitem{Strassler1992}
M.~J. Strassler, ``Field theory without Feynman diagrams: One-loop effective actions,''
Nucl. Phys. B \textbf{385}, 145--184 (1992).

\bibitem{Feynman1950}
R.~P. Feynman, ``Mathematical formulation of the quantum theory of electromagnetic interaction,''
Phys. Rev. \textbf{80}, 440--457 (1950).

\bibitem{DuttonGWFM}
J.~D. Dutton, \textit{Geometric Worldline Field Model (GWFM)}, unpublished manuscript (2026).

\bibitem{DuttonSEFI}
J.~D. Dutton, \textit{Single Entity Field Interpretation (SEFI) Canon}, unpublished manuscript (2026).

\end{thebibliography}

\end{document}
